\section{AION Flow Simulation, Governed Execution, Recovery and Route Receipts}
\label{sec:aion-flow-stage-eight}

\subsection{Purpose}

Stage Eight converts the sovereign visual intelligence graph into a controlled runtime
without allowing a model node to become an execution authority.  The existing Boardroom Workflow
Canvas remains the authoring surface.  A new \textbf{Run \& receipts} control exposes compilation,
simulation, exact review, execution state, recovery controls and the final route receipt in the same
interface.

The operating sequence is:

\begin{quote}
Visible graph $\rightarrow$ sovereign compile $\rightarrow$ Workflow Capsule $\rightarrow$
static validation $\rightarrow$ policy and disclosure preflight $\rightarrow$ simulation
$\rightarrow$ exact review $\rightarrow$ authorized capability execution $\rightarrow$
verification $\rightarrow$ Intelligence Route Receipt.
\end{quote}

Drawing, saving, compiling or simulating a workflow causes no external write.  Model output remains
a proposal and cannot call a connector directly.

\subsection{Compilation and Workflow Capsule Translation}

The runtime normalizes nodes and edges from the visible Glyph canvas and submits the graph to the
Sovereign Flow compiler.  Compilation requires exactly one AION admission boundary and one AION
receipt boundary.  Every node must be reachable from admission and able to reach the receipt.
Capabilities require an upstream policy node, while consequential capabilities additionally require
an upstream approval node.  External destinations must declare their provider, data classification,
destination, residency and encrypted edge.

A successful sovereign manifest is translated into the existing
\texttt{aion.workflow\_capsule.v1} contract.  This preserves the deployed Workflow Capsule system
rather than creating a parallel executor.  The capsule records its sovereign manifest hash, ordered
steps, graph links, dry-run-first policy and approval-before-write rule.  Compilation explicitly marks
the result as not executed.

\subsection{Static Validation and Safe Simulation}

Before a run can be authorized, Stage Eight checks:

\begin{itemize}
  \item node and edge identities and supported node types;
  \item capability and verification bindings;
  \item the AION admission and receipt boundary;
  \item policy and approval placement;
  \item external disclosure, encryption and residency declarations; and
  \item graph cycles, which are rejected unless expressed through a separately bounded deliberation
        node with explicit limits.
\end{itemize}

The simulator produces sample output hashes and per-node simulated states.  It performs zero
external writes.  Its report lists predicted boundary crossings, exact fields and destinations,
estimated cost and the number of consequential capabilities.  Secrets and secret-shaped fields are
redacted before any preparation or run record is persisted.

\subsection{Exact Review and Approval}

The preparation generates a canonical review containing the flow identifier, manifest hash, active
person and workspace, capability list, targets, consequential-action flags, disclosure-report hash,
estimated cost and approval count.  A canonical \texttt{review\_hash} binds these facts together.

The user must open \textbf{Review exact route}, confirm the displayed route and explicitly authorize
it.  If approval changes the authorized state, AION generates a new review hash and requires final
confirmation before execution.  A changed manifest, expired authority, revoked identity, missing
approval or separation-of-duty violation fails closed.  Shared model output cannot manufacture this
approval.

\subsection{Execution Boundary}

Non-consequential local capabilities may execute through the bounded local capability interface.
Consequential capabilities are dispatched only through the existing permissioned Workflow Connector
adapter.  The default runtime refuses an unbound consequential capability.  Connector credentials
remain opaque mother-brain vault references and never enter the graph, browser state or receipt.

Connector readiness and external completion are different facts.  A route may prove that its policy,
approval and connector prerequisites are ready without claiming that an email, record or other
external action occurred.  The final receipt marks an external effect as verified only when the
permissioned adapter returns a positive provider outcome for that exact operation.

\subsection{Run State, Idempotency and Recovery}

Each run is journalled atomically in the customer-owned runtime directory.  Every node has an ordered
state selected from queued, running, waiting, blocked, failed and verified.  The journal records
attempts and events without storing raw credentials.

Capability execution uses a stable idempotency key.  A resumed operation that has already been
verified is not issued again.  Retry counts and timeouts are bounded, and an open connector circuit
stops further attempts.  Compensation instructions are executed automatically only when explicitly
marked safe and local; otherwise they enter \texttt{waiting\_approval} rather than being falsely
reported as complete.

The Boardroom exposes \textbf{Pause}, \textbf{Cancel}, \textbf{Resume} and \textbf{Recover}.  After a
mother-brain restart, a run found in the running state is moved to waiting and its interrupted node is
returned to queued.  Resume requires the same stored manifest hash.  Completed, cancelled or
compensated runs cannot be silently restarted.

\subsection{Intelligence Route Receipt}

Every terminal attempt produces \texttt{aion.intelligence\_route\_receipt.v1}.  The receipt contains:

\begin{itemize}
  \item receipt, run and flow identifiers;
  \item sovereign manifest and exact-review hashes;
  \item final run status;
  \item each node's final state and output hash;
  \item whether all consequential external effects were actually provider-verified; and
  \item issue time and a canonical tamper-evident receipt hash.
\end{itemize}

The receipt is visible inside \textbf{Run \& receipts} and can be inspected without exposing the
underlying restricted content.  It proves the route that was reviewed and the outcome that was
observed; it does not turn model agreement or connector readiness into evidence of execution.

\subsection{Boardroom Location and Use}

The capability is deliberately visible in the existing Boardroom rather than hidden in a developer
API.  Open the Boardroom's \textbf{Workflow Canvas}.  Beside \textbf{Execute workflow}, select
\textbf{Run \& receipts}, then:

\begin{enumerate}
  \item select \textbf{Compile \& simulate};
  \item inspect static validation, policy, predicted cost, disclosure and simulated node states;
  \item select \textbf{Review exact route};
  \item confirm and authorize the exact route when an approval is required;
  \item confirm the final review hash and select \textbf{Run authorized route};
  \item monitor node states or use pause, cancel, resume and recover; and
  \item inspect the resulting \textbf{Intelligence Route Receipt}.
\end{enumerate}

This surface is intended for ordinary Boardroom operators.  Advanced payloads remain inspectable but
are not required to understand whether a route is safe, blocked, waiting, complete or externally
verified.

\subsection{Security Principle}

\begin{quote}
\textbf{AION may compile and coordinate the intelligence route, but only customer policy, exact human
authority and a permissioned capability adapter can create a consequential external effect.}
\end{quote}
